%=============================================================================
\section{Digital Signal Processing}
%=============================================================================

You might be wondering: ``I signed up for control theory, why are we starting with signal processing?'' Great question. Here is the answer: \textbf{garbage in, garbage out.}

Your controller is only as good as the data it receives. If your IMU (Inertial Measurement Unit---the small chip that combines accelerometers and gyroscopes to measure orientation and motion) reading is bouncing around like a caffeinated squirrel, no amount of PID tuning will save you. So before we learn to control systems, we learn to clean up their signals.

\subsection{Understanding the Frequency Domain}

In everyday life, we think about signals in the \textbf{time domain}---voltage changes over time, position changes over time. This is natural and intuitive. But it turns out that many problems become \textit{much} easier if we look at signals from a different angle: the \textbf{frequency domain}.

The core idea is beautifully simple: \textit{any signal can be decomposed into a sum of sinusoids of different frequencies.} Your noisy gyroscope reading? It is actually a clean low-frequency rotation signal plus a bunch of high-frequency noise signals all added together. If we can identify which frequencies are ``signal'' and which are ``noise,'' we can surgically remove the noise.

Think of it like an audio equalizer on a music player. The bass, mids, and treble are all mixed together in the time-domain waveform. But the equalizer lets you boost or cut specific frequency ranges independently. That is exactly what we do with filters in control systems.

\textbf{Key intuition:} Slow, meaningful changes in your system (like a chassis turning) happen at low frequencies. Fast, meaningless changes (like electrical noise or vibration) happen at high frequencies. The frequency domain lets us tell them apart.

\subsection{Laplace Transform and Fourier Transform}

These two transforms are your tickets from the time domain to the frequency domain. They are closely related, but they serve slightly different purposes.

\subsubsection{Laplace Transform}

The Laplace transform takes a time-domain function $f(t)$ and converts it into a function of the complex variable $s$:

\begin{equation}
F(s) = \mathcal{L}\{f(t)\} = \int_0^{\infty} f(t) e^{-st} \, dt
\end{equation}

where $s = \sigma + j\omega$ is a complex number. The real part $\sigma$ captures exponential growth or decay, and the imaginary part $j\omega$ captures oscillation.

\textbf{Why do we care?} Because the Laplace transform turns differential equations into algebraic equations. Instead of solving
\[
m\ddot{x} + b\dot{x} + kx = F(t),
\]
(which requires calculus), we can solve
\[
(ms^2 + bs + k)X(s) = F(s),
\]
(which is just algebra). Derivatives become multiplication by $s$. Integration becomes division by $s$. Life becomes easier.

The most important Laplace transforms you will use daily:

\begin{center}
\begin{tabular}{lll}
\toprule
\textbf{Time domain} & \textbf{$s$-domain} & \textbf{When you see it} \\
\midrule
$1$ (step input) & $\frac{1}{s}$ & Commanding a new setpoint \\
$e^{-at}$ & $\frac{1}{s+a}$ & First-order system response \\
$\sin(\omega t)$ & $\frac{\omega}{s^2 + \omega^2}$ & Oscillation, vibration \\
$\dot{f}(t)$ & $sF(s) - f(0)$ & Derivative in the controller \\
$\int f(t)\,dt$ & $\frac{F(s)}{s}$ & Integral in the controller \\
\bottomrule
\end{tabular}
\end{center}

\textbf{Pro tip:} You do not need to memorize the integral formula. You need to understand that $s$ means ``derivative'' and $1/s$ means ``integral.'' That single insight will carry you through 80\% of control theory.

\textbf{What is $e^{-st}$ actually doing?} The kernel $e^{-st} = e^{-\sigma t} e^{-j\omega t}$ is a ``probe signal'' that asks two questions simultaneously: (1) $e^{-\sigma t}$: ``Does the signal grow or decay at rate $\sigma$?'' (2) $e^{-j\omega t}$: ``How much of frequency $\omega$ is present?'' By integrating $f(t) \cdot e^{-st}$ over all time, we compute the ``correlation'' between the signal and every possible exponential-sinusoid. The result $F(s)$ is a map of the entire complex plane, telling us exactly how much of each exponential-oscillatory mode exists in the signal.

\textbf{Common Laplace transform pairs --- worked out:}

The \textbf{unit step} $u(t) = 1$ for $t \geq 0$:
\[
\mathcal{L}\{1\} = \int_0^{\infty} e^{-st} \, dt = \left[-\frac{e^{-st}}{s}\right]_0^{\infty} = \frac{1}{s}
\]

The \textbf{exponential decay} $e^{-at}$:
\[
\mathcal{L}\{e^{-at}\} = \int_0^{\infty} e^{-at} e^{-st} \, dt = \int_0^{\infty} e^{-(s+a)t} \, dt = \frac{1}{s+a}
\]
Notice that $F(s) = \frac{1}{s+a}$ ``blows up'' when $s = -a$. That special value of $s$ is called a \textbf{pole}---we will define it properly in a moment.

\subsubsection{A First Look at Poles and Zeros}

We just saw that $\mathcal{L}\{e^{-at}\} = \frac{1}{s+a}$. The denominator becomes zero at $s = -a$, which means $F(s)$ shoots to infinity there. That point is called a \textbf{pole}.

In general, for any rational function in $s$---and most \textbf{transfer functions} (the ratio of a system's output to its input in the $s$-domain, which we will derive properly in the next chapter) are rational---we write:
\[
G(s) = \frac{N(s)}{D(s)} = \frac{(s - z_1)(s - z_2) \cdots}{(s - p_1)(s - p_2) \cdots}
\]
\begin{itemize}[nosep]
    \item \textbf{Poles} ($p_i$): roots of the denominator $D(s)$. These are the values of $s$ where $G(s) \to \infty$.
    \item \textbf{Zeros} ($z_i$): roots of the numerator $N(s)$. These are the values of $s$ where $G(s) = 0$.
\end{itemize}

\textbf{Why should you care right now?} Because the pole location tells you the time-domain behavior directly:
\begin{itemize}[nosep]
    \item A pole at $s = -a$ (negative real) produces $e^{-at}$---exponential decay. The more negative $a$ is, the faster the decay.
    \item A pole at $s = +a$ (positive real) produces $e^{+at}$---exponential growth. The system is unstable.
    \item A pair of complex poles $s = -\sigma \pm j\omega$ produces $e^{-\sigma t}\cos(\omega t)$---a decaying oscillation. (This should look familiar from the $e$ and $\pi$ section!)
\end{itemize}

In short: \textbf{poles tell you what exponential and oscillatory modes are hiding inside a system.} Every time we design a filter in this chapter, we are really choosing where to put poles and zeros. We will revisit poles and zeros more rigorously in Section~\ref{sec:system_analysis_tf}, where we connect them to stability, speed, and overall system performance.

The \textbf{ramp} $t \cdot u(t)$:
\[
\mathcal{L}\{t\} = \frac{1}{s^2}
\]
Notice: multiplying by $t$ in time corresponds to differentiating with respect to $s$ and flipping the sign: $\mathcal{L}\{t \cdot f(t)\} = -\frac{d}{ds}F(s)$. Since $\mathcal{L}\{1\} = 1/s$, we get $\mathcal{L}\{t\} = -\frac{d}{ds}\frac{1}{s} = \frac{1}{s^2}$.

\subsubsection{Fourier Transform}

The Fourier transform is the Laplace transform's well-behaved cousin. It is what you get when you set $s = j\omega$ (i.e., $\sigma = 0$):

\begin{equation}
F(j\omega) = \int_{-\infty}^{\infty} f(t) e^{-j\omega t} \, dt
\end{equation}

While the Laplace transform works with the full complex plane (and can handle unstable, growing signals), the Fourier transform lives strictly on the imaginary axis. It tells you: ``How much of each frequency is present in this signal?''

\textbf{When to use which:} Use the \textbf{Laplace transform} for analyzing and designing control systems (transfer functions, stability). Use the \textbf{Fourier transform} for analyzing signals (noise, vibration, frequency content).

In practice, for digital systems we often use the \textbf{Discrete Fourier Transform (DFT)}, computed efficiently via the \textbf{Fast Fourier Transform (FFT)} algorithm. If you ever want to see what frequencies are hiding in your sensor data, just throw it into an FFT. MATLAB's \texttt{fft()} or Python's \texttt{numpy.fft.fft()} will do the job in one line.

\subsubsection{Frequency Spectrum}

When you apply the Fourier transform to a signal, you get its \textbf{frequency spectrum}---a plot showing the amplitude (and phase) of each frequency component.

Reading a frequency spectrum is like reading an X-ray of your signal:
\begin{itemize}
    \item A tall peak at 50 Hz? That is probably power line interference.
    \item A broad hump below 10 Hz? That is likely your actual signal of interest.
    \item Random fuzz spread across all frequencies? That is white noise.
\end{itemize}

\textbf{Practical example:} Suppose your gimbal encoder gives a noisy velocity reading. You record 1 second of data, compute the FFT, and see a clean peak at 3 Hz (your gimbal is tracking a target moving at 3 Hz) plus a mess of noise above 50 Hz. Now you know: design a low-pass filter with a cutoff around 20--30 Hz, and your signal will be clean.

The frequency spectrum has two parts:
\begin{itemize}
    \item \textbf{Magnitude spectrum} $|F(j\omega)|$: How strong each frequency component is.
    \item \textbf{Phase spectrum} $\angle F(j\omega)$: The time shift of each frequency component.
\end{itemize}

For filter design, we mostly care about the magnitude spectrum. Phase becomes important later when we analyze control loop stability (\textit{phase margin}---how much extra phase lag the system can tolerate before it becomes unstable; see the Classical Control chapter).

\begin{figure}[H]
\centering
\includegraphics[width=\textwidth]{frequency_spectrum.pdf}
\caption{Top: a noisy time-domain signal looks hopeless. Bottom: the FFT instantly reveals a 5 Hz signal and a 42 Hz resonance hiding in the noise.}
\end{figure}

\vspace{0.5em}
\begin{mdframed}[linewidth=1pt, innertopmargin=10pt, innerbottommargin=10pt]
\textbf{Case Study: Debugging a Vibrating Chassis with FFT}

\vspace{0.3em}
True story pattern from competition robotics: Your mecanum chassis drives fine at low speed but develops a nasty vibration at high speed. The whole frame buzzes. Everyone suspects a loose screw. You tighten everything. Still vibrates.

\textbf{The diagnosis:} Log the IMU accelerometer data for 5 seconds while driving at the problematic speed. Run an FFT in Python:

\begin{verbatim}
    import numpy as np
    freqs = np.fft.rfftfreq(len(data), d=1.0/sample_rate)
    spectrum = np.abs(np.fft.rfft(data))
    plt.plot(freqs, spectrum)
\end{verbatim}

You see a massive spike at exactly 42 Hz. Interesting. You calculate: at high speed, the mecanum wheel rollers contact the ground at\ldots about 42 Hz. The roller contact frequency matches a structural resonance of the chassis frame.

\textbf{The fix options:}
\begin{itemize}[nosep]
    \item Add a \textbf{notch filter} (band-reject) at 42 Hz in the control loop to prevent the controller from exciting the resonance.
    \item Add physical damping (rubber mounts) to shift the resonance frequency.
    \item Stiffen the frame to move the resonance above the operating range.
\end{itemize}

\textbf{The lesson:} Without the FFT, you would have spent hours tightening screws. With the FFT, the root cause was obvious in 10 seconds. \textbf{The frequency domain is a debugging superpower.} Whenever something oscillates and you do not know why, record the data and look at the spectrum.
\end{mdframed}

\subsection{Frequency-Based Filter Design}

Now that we can see signals in the frequency domain, we can design filters to keep the frequencies we want and reject the ones we do not.

A filter is simply a system that has different gains at different frequencies. We describe this using the filter's \textbf{frequency response}---a plot of gain vs.\ frequency (called a \textbf{Bode magnitude plot}).

\subsubsection{Low-Pass Filter}

The most common filter in robotics. It passes low frequencies and attenuates high frequencies.

\textbf{First-order low-pass filter:}
\begin{equation}
H(s) = \frac{\omega_c}{s + \omega_c}
\end{equation}

where $\omega_c = 2\pi f_c$ is the cutoff frequency in rad/s. At frequencies below $\omega_c$, the gain is approximately 1 (signal passes through). Above $\omega_c$, the gain rolls off at $-20$ dB/decade (signal gets attenuated).

\textbf{In the time domain}, this is equivalent to an exponential moving average:
\begin{equation}
y[n] = \alpha \cdot x[n] + (1 - \alpha) \cdot y[n-1]
\end{equation}

This follows from a backward Euler discretization of the continuous filter. Starting from $H(s) = \frac{\omega_c}{s + \omega_c}$, the corresponding differential equation is $\dot{y} + \omega_c y = \omega_c x$. Replacing $\dot{y}$ with the backward difference $\frac{y[n] - y[n-1]}{T_s}$ gives:
\[
\frac{y[n] - y[n-1]}{T_s} + \omega_c\, y[n] = \omega_c\, x[n]
\]
Collecting terms in $y[n]$:
\[
y[n]\!\left(\frac{1}{T_s} + \omega_c\right) = \omega_c\, x[n] + \frac{y[n-1]}{T_s}
\]
Dividing both sides by $\left(\frac{1}{T_s} + \omega_c\right)$ and defining $\alpha = \frac{\omega_c T_s}{1 + \omega_c T_s} = \frac{T_s}{T_s + 1/\omega_c}$:
\[
y[n] = \alpha \cdot x[n] + (1 - \alpha) \cdot y[n-1]
\]

where $\alpha = \frac{T_s}{T_s + \frac{1}{2\pi f_c}}$ (since $\omega_c = 2\pi f_c$) and $T_s$ is the sampling period.

\textbf{When to use:} Smoothing noisy sensor readings (gyroscope, accelerometer, motor current). This is the filter you will use 90\% of the time.

\textbf{Gotcha:} A low-pass filter introduces \textbf{phase lag}---the output is delayed relative to the input. The more aggressively you filter (lower cutoff), the more lag you introduce. In a fast control loop, this lag can destabilize your system. There is always a trade-off between smooth signals and fast response.

\begin{figure}[H]
\centering
\includegraphics[width=\textwidth]{lpf_response.pdf}
\caption{Left: LPF step response --- lower cutoff means slower rise but smoother output. Right: filtering a noisy 3 Hz signal --- the LPF recovers the true signal while removing high-frequency noise.}
\end{figure}

\vspace{0.5em}
\begin{mdframed}[linewidth=1pt, innertopmargin=10pt, innerbottommargin=10pt]
\textbf{Case Study: Why a Low-Pass Filter Gives You ``Soft Start'' for Free}

\vspace{0.3em}
Suppose you command your motor to jump from 0 to 5000 RPM instantly (a step input). Without any filter, the setpoint changes discontinuously---the motor driver sees a sudden full-voltage command, the gears slam, the chassis lurches, and your teammates spill their coffee.

Now insert a first-order LPF on the setpoint: $r_{\text{filtered}}(s) = \frac{\omega_c}{s + \omega_c} \cdot R(s)$. What happens?

The step input in the $s$-domain is $R(s) = \frac{5000}{s}$. After the filter:
\[
R_{\text{filtered}}(s) = \frac{5000 \, \omega_c}{s(s + \omega_c)}
\]

Taking the inverse Laplace transform:
\[
r_{\text{filtered}}(t) = 5000 \left(1 - e^{-\omega_c t}\right)
\]

The step has been transformed into a \textbf{smooth exponential ramp}. The time constant is $\tau = 1/\omega_c$. After $3\tau$, the output reaches 95\% of the target; after $5\tau$, it is at 99\%.

\textbf{Why this matters on your robot:}
\begin{itemize}[nosep]
    \item \textbf{Mechanical protection:} Gears and belts have torque limits. A step command can strip teeth or snap belts. The LPF limits the rate of change, keeping forces within safe bounds.
    \item \textbf{Current limiting:} A sudden voltage step causes a huge inrush current ($I = V/R$ when back-EMF is still zero). The LPF ramps up the voltage gradually, keeping current manageable.
    \item \textbf{Tunable aggressiveness:} Want a faster ramp? Increase $\omega_c$ (higher cutoff). Want a gentler ramp? Decrease $\omega_c$. One parameter controls the entire behavior.
\end{itemize}

\textbf{Comparison with a linear ramp:} You could also ramp the setpoint linearly (e.g., increase by 1000 RPM/s). But the LPF is better because it also smooths out \textit{any} discontinuity in the setpoint---not just the initial step, but also sudden changes during operation. It is a ``set it and forget it'' solution.

\textbf{The code is trivial:}
\begin{verbatim}
    setpoint_filtered += alpha * (setpoint_raw - setpoint_filtered);
\end{verbatim}

One line of code. No lookup tables, no state machines. That is the beauty of the LPF as a soft-start mechanism.
\end{mdframed}

\begin{figure}[H]
\centering
\includegraphics[width=\textwidth]{lpf_soft_start.pdf}
\caption{LPF soft start: a single parameter $\tau$ controls how aggressively the motor ramps up. The step command becomes a smooth exponential curve.}
\end{figure}

\subsubsection{Band-Pass Filter}

Passes a specific range of frequencies and attenuates everything else. Think of it as a low-pass and a high-pass filter combined.

\begin{equation}
H(s) = \frac{\frac{\omega_0}{Q} s}{s^2 + \frac{\omega_0}{Q} s + \omega_0^2}
\end{equation}

where $\omega_0$ is the center frequency and $Q$ is the quality factor (higher $Q$ = narrower band).

\textbf{When to use:} Extracting a specific frequency component from a signal. For example, if you know your target oscillates at a specific frequency and you want to track it while rejecting everything else.

In robotics, band-pass filters are less common than low-pass filters, but they are useful for vibration analysis and resonance detection.

\subsubsection{High-Pass Filter}

Passes high frequencies, attenuates low frequencies. It is the opposite of a low-pass filter.

\textbf{First-order high-pass filter:}
\begin{equation}
H(s) = \frac{s}{s + \omega_c}
\end{equation}

\textbf{When to use:} Removing DC offset or slow drift from a signal. For example, if your accelerometer has a slowly drifting bias, a high-pass filter can remove it while keeping the fast acceleration data intact.

\textbf{Practical note:} A complementary filter (commonly used for IMU sensor fusion) is just a low-pass filter on one sensor and a high-pass filter on another, with their cutoff frequencies matched so the gains add up to 1. The accelerometer gets low-passed (trust it for slow/static orientation) and the gyroscope gets high-passed (trust it for fast rotations). Simple, elegant, and surprisingly effective.

\subsubsection{Notch Filter (Band-Reject Filter)}

A notch filter is the opposite of a band-pass filter: it \emph{removes} a narrow band of frequencies while leaving everything else untouched. It is your surgical tool for killing a single problematic frequency---a structural resonance, a 50/60 Hz power-line hum, or a motor cogging frequency.

\textbf{Second-order notch filter:}
\begin{equation}
H(s) = \frac{s^2 + \omega_0^2}{s^2 + \frac{\omega_0}{Q}\,s + \omega_0^2}
\end{equation}

where $\omega_0$ is the frequency to reject and $Q$ is the quality factor (higher $Q$ = narrower notch). At $s = j\omega_0$, the numerator evaluates to zero, so the gain drops to zero at exactly $\omega_0$. The denominator keeps the filter stable and controls the width of the rejection band.

\textbf{Pole-zero interpretation:} The numerator places a pair of \textit{zeros} on the imaginary axis at $s = \pm j\omega_0$---these create the ``notch'' by driving the gain to zero. The denominator places a pair of \textit{poles} slightly to the left of those zeros---these keep the filter stable and determine how quickly the gain recovers on either side. This is a textbook example of poles and zeros doing exactly what we described in Section 2.2: zeros kill frequencies, poles shape the surrounding response.

\textbf{When to use:}
\begin{itemize}[nosep]
    \item A structural resonance is causing oscillation in your control loop (see the mecanum wheel case study earlier in this chapter).
    \item Power-line interference (50 or 60 Hz) is corrupting a sensor signal.
    \item A motor's cogging torque at a known frequency is visible in the velocity measurement.
\end{itemize}

\textbf{Practical tip:} Getting $\omega_0$ wrong by even a few Hz can make the notch miss the target entirely. Always identify the exact resonance frequency from an FFT of your sensor data \textit{before} designing the filter. If the resonance frequency shifts with operating conditions (e.g., with speed or load), consider an \textbf{adaptive notch filter} that tracks $\omega_0$ in real time.

\textbf{Discrete-time implementation:} Like all filters in this chapter, the notch filter is implemented as a biquad (second-order section). The bilinear transform with pre-warping at $\omega_0$ converts $H(s)$ to $H(z)$ without shifting the notch frequency:
\[
s = \frac{\omega_0}{\tan(\omega_0 T_s / 2)} \cdot \frac{z - 1}{z + 1}.
\]

\subsubsection{Filter Type Selection and Parameter Tuning}

Choosing the right filter comes down to answering three questions:
\begin{enumerate}
    \item \textbf{What frequencies do I want to keep?} This determines the filter type (low-pass, high-pass, band-pass, or notch).
    \item \textbf{Where is the boundary?} This determines the cutoff frequency $f_c$ (or the rejection frequency $f_0$ for a notch).
    \item \textbf{How sharp should the transition be?} This determines the filter order (or $Q$ for a notch).
\end{enumerate}

\textbf{Rules of thumb:}
\begin{itemize}
    \item Start with a first-order low-pass filter. It is simple, stable, and good enough for most applications.
    \item Set the cutoff frequency to 5--10$\times$ the bandwidth of your signal of interest. Too low and you lose signal; too high and you keep noise.
    \item If a first-order filter is not enough, try second-order (Butterworth for flat passband, Chebyshev for sharper rolloff).
    \item \textbf{Never} design a filter without looking at the frequency spectrum of your actual data first. Guessing cutoff frequencies is a recipe for frustration.
    \item When in doubt, plot the filtered output alongside the raw signal and check if the result makes physical sense.
\end{itemize}

\subsection{Z-Transform}

Everything we have discussed so far with the Laplace transform lives in the continuous world. But your microcontroller is digital---it works with samples, not continuous signals. The \textbf{Z-transform} is the discrete-time equivalent of the Laplace transform.

The Z-transform of a discrete sequence $x[n]$ is:
\begin{equation}
X(z) = \sum_{n=0}^{\infty} x[n] z^{-n}
\end{equation}

The key correspondence is:
\begin{center}
\begin{tabular}{ll}
\toprule
\textbf{Continuous ($s$-domain)} & \textbf{Discrete ($z$-domain)} \\
\midrule
$s$ (derivative) & $\frac{z-1}{T_s}$ (forward difference) \\
$\frac{1}{s}$ (integral) & $\frac{T_s}{z-1}$ (accumulator) \\
Stable if $\text{Re}(s) < 0$ & Stable if $|z| < 1$ \\
$e^{-aT_s}$ in $s$-domain & $z^{-1}$ is one-sample delay \\
\bottomrule
\end{tabular}
\end{center}

\textbf{Why do we care?} Because when you implement a filter or controller on a microcontroller, you are working in the $z$-domain, whether you realize it or not. That simple line of code:
\[
\texttt{y = alpha * x + (1 - alpha) * y\_prev;}
\]
is actually a $z$-domain transfer function:
\[
H(z) = \frac{\alpha}{1 - (1-\alpha)z^{-1}}
\]

Understanding the Z-transform helps you:
\begin{itemize}
    \item Convert continuous-time filters/controllers to code (covered in Section 5).
    \item Analyze the stability of your discrete-time system.
    \item Understand why your controller might behave differently at different sampling rates.
\end{itemize}

For now, just remember: \textbf{$z^{-1}$ means ``one time step ago.''} That is the single most important thing about the Z-transform.

\textbf{Common Z-transform pairs:}

\begin{center}
\begin{tabular}{lll}
\toprule
\textbf{Discrete signal} & \textbf{$z$-domain} & \textbf{When you see it} \\
\midrule
$\delta[n]$ (unit impulse) & $1$ & Initial kick \\
$u[n]$ (unit step) & $\frac{z}{z-1}$ & Step command \\
$a^n u[n]$ (geometric decay) & $\frac{z}{z-a}$ & Discrete first-order response \\
$x[n-1]$ (one-step delay) & $z^{-1}X(z)$ & Previous sample in code \\
$x[n] - x[n-1]$ (first difference) & $(1 - z^{-1})X(z)$ & Discrete derivative \\
$\sum_{k=0}^{n} x[k]$ (running sum) & $\frac{1}{1 - z^{-1}}X(z)$ & Discrete integral \\
\bottomrule
\end{tabular}
\end{center}

\textbf{The connection between $s$ and $z$:} The exact mapping is $z = e^{sT_s}$. The connection is simple: a continuous-time exponential $x(t) = e^{st}$, when sampled at intervals $T_s$, becomes the sequence $x[n] = e^{snT_s} = (e^{sT_s})^n$. In the Z-transform world, $z^n$ is the discrete exponential. Matching terms gives $z = e^{sT_s}$.

A continuous pole at $s = -a$ maps to a discrete pole at $z = e^{-aT_s}$. This means:
\begin{itemize}
    \item A stable continuous pole ($\text{Re}(s) < 0$) maps to $|z| < 1$ (inside the unit circle). Stability is preserved.
    \item A faster continuous pole (more negative $s$) maps to a smaller $|z|$ (closer to the origin). Very fast poles map near $z = 0$.
    \item An oscillatory continuous pole ($s = \sigma \pm j\omega$) maps to $z = e^{\sigma T_s} e^{\pm j\omega T_s}$, which spirals inside (stable) or outside (unstable) the unit circle.
\end{itemize}

\textbf{Reading a $z$-domain transfer function:} Any $z$-domain transfer function can be directly translated to code by replacing $z^{-1}$ with ``previous value.'' For example:
\[
H(z) = \frac{b_0 + b_1 z^{-1}}{1 + a_1 z^{-1}}
\]
becomes:
\begin{verbatim}
    y = b0 * x + b1 * x_prev - a1 * y_prev;
    x_prev = x;  y_prev = y;
\end{verbatim}
This is called a \textbf{difference equation}, and it is the final form of any digital filter or controller.
